% A PREAMBLE IS READ IN THE ALIGNMENT'S OWN MODE, AND ITS TEMPLATES ARE
% NOT.
%
% The preamble of an \halign is read in internal vertical mode and of a
% \valign in restricted horizontal mode -- the alignment's own list. But
% almost nothing in a preamble is expanded as it is read: the templates
% are kept as they were written and run later, inside each entry, in the
% entry's own mode. \span is the exception, and forces the token after it
% to be expanded where it stands:
%
%   \halign{&#\span\iftrue...}   {internal vertical mode: \iftrue}
%   \halign{\iftrue#\fi...}      {restricted horizontal mode: \iftrue}
%
% both for the same \iftrue written in the same preamble.
%
% AND `end of alignment template' comes BEHIND the v part, not in front
% of it. The reference reads a frozen \endtemplate once the v part is
% exhausted and traces it there, so
%
%   \halign{#\span\A Y&#\cr A&B\cr}     with \def\A{X}
%
% draws {the letter A}, {the letter X}, {the letter Y} and only then
% {end of alignment template}. HSTeX drew it as the entry's own text
% ended, in front of the whole v part.
%
% trip line 331 writes `\halign spread-12.truedd{&#\span\iftrue\A\span
% \else\span\fi\span&'.
\catcode`\{=1 \catcode`\}=2 \catcode`\#=6 \catcode`\&=4
\tracingonline=1 \hsize=100pt \parindent=0pt \tabskip=0pt
\hbadness=10000 \vbadness=10000 \hfuzz=1000pt \vfuzz=1000pt
\def\A{X}
\message{[A span forces expansion in a halign preamble]}
\tracingcommands=2
\halign{&#\span\iftrue\A\span\else\span\fi\span\cr A\cr}
\tracingcommands=0
\message{[B and in a valign preamble]}
\tracingcommands=2
\noindent P\valign{&#\span\iftrue\A\span\else\span\fi\span\cr A\cr}Q\par
\tracingcommands=0
\message{[C the v part runs before the template ends]}
\tracingcommands=2
\halign{#\span\A Y&#\cr A&B\cr}
\tracingcommands=0
\message{[D a plain expandable in a preamble runs with the entry]}
\tracingcommands=2
\halign{\iftrue#\fi Y&#\cr A&B\cr}
\tracingcommands=0
\message{[done]}
\end
